\paragraph{Explain the insurer's annual operating cash from the same journals.}
The earlier indirect bridge can now use each fiscal report directly. Under one
supplied illustrative policy, liability movements from finance are noncash
profit adjustments, liability releases are working-capital adjustments, and
all six methods' income and expense postings remain in operating profit.
These choices are premises of the explanation. They do not alter the native
journals or the operating classification already supplied for claim payments.

Using the earlier meanings of $P$, $W$, $N$ and $R$, the two nonempty years give
\[
\begin{aligned}
  2026:\quad C_o&=-5-55+10+0=-50,\\
  2027:\quad C_o&=10-50-5+0=-45.
\end{aligned}
\]
The negative working-capital terms reflect the liability releases; the finance
terms reverse the corresponding noncash expense or income. The selected 2025
report has no events and all terms are zero. The loss component $K$ contributes
no extra monetary adjustment because it is memory within the liability,
whose postings are already included. Adding $K$ again would count part of the
same change twice.

A second supplied policy classifies all liability movements as noncash profit
adjustments. It gives $W=0$ and $N=-45$ in 2026, then $W=0$ and $N=-55$ in
2027, with the same profit, cash and financial reports. Exact reconciliation
therefore leaves this presentation choice unresolved. Both policies preserve
every gross posting, including the zero loss reversal in the zero-loss case.
Removing its classification is an error even though its amount is zero.

The comparison selects each period after its preceding carry and binds the
exact fiscal method hashes. A pre-carry method hash or a classification from
another year cannot stand in for the required policy. Complete exchanges under
\path{data/accounting-master/insurance-statement-bridge-20261025/} cover the
two charts, two retained native cases, two explicit policies and three years,
with the unrelated equipment balance present. They use the existing indirect
engine and reproduce the original annual financial reports. Actual
classification validity and event completeness remain separate questions.
